\documentclass[a4paper,12pt]{extarticle}

% ============================================================
%                    ЯЗЫК И ШРИФТЫ
% ============================================================

% Компиляция: XeLaTeX
\usepackage{fontspec}
\usepackage{polyglossia}

\setdefaultlanguage{russian}

\setmainfont{Times New Roman}[Script=Cyrillic,Ligatures=TeX]
\newfontfamily\cyrillicfont{Times New Roman}[Script=Cyrillic,Ligatures=TeX]
\newfontfamily\calibrifont{Liberation Sans}
\newfontfamily\cyrillicfonttt{Courier New}



% ============================================================
%                    МАТЕМАТИКА
% ============================================================

\usepackage{amsmath,amssymb}


% ============================================================
%                    ЦВЕТА И ГРАФИКА
% ============================================================

\usepackage[table]{xcolor}
\usepackage{graphicx}
\usepackage{pgfplots}
\usepackage{float}
\usepackage{multicol}


% ============================================================
%                    ТАБЛИЦЫ
% ============================================================

\usepackage{tabularx}
\usepackage{array}
\usepackage{booktabs}
\usepackage{multirow}
\usepackage{makecell}
\usepackage{tabularray}
\usepackage{ragged2e}
\usepackage{longtable}

% ============================================================
%                    СПИСКИ
% ============================================================

\usepackage{enumitem}


% ============================================================
%                    URL И ССЫЛКИ
% ============================================================

\usepackage{xurl}



% ============================================================
%                    ГИПЕРССЫЛКИ
% ============================================================

\usepackage{xurl}
\usepackage{microtype}

\usepackage[
    left=30mm,
    right=15mm,
    top=20mm,
    bottom=20mm
]{geometry}



\begin{document}

\begin{titlepage}
    \centering


    {\LARGE\bfseries Задача A1\par}
    \vspace{0.4cm}
    {\Large\bfseries Анализ строковых сортировок\par}


    \vspace{0.5cm}

    \begin{center}
    \begin{tabular}{|c|c|}
        \hline
        \textbf{Задача} & \textbf{ID посылки} \\
        \hline
        A1m & \texttt{375963079} \\
        \hline
        A1q & \texttt{375963355} \\
        \hline
        A1r & \texttt{375966130} \\
        \hline
        A1rq & \texttt{375969103} \\
        \hline
    \end{tabular}
    \end{center}

    \vspace{16cm}

    \begin{flushright}
        \large
        \textbf{Выполнил:} Коробов Вячеслав\\
        \textbf{Группа:} 246
    \end{flushright}



\end{titlepage}

\section{Эмпирический анализ стандартных алгоритмов}

Для стандартных алгоритмов \texttt{QUICKSORT} и \texttt{MERGESORT} были измерены две характеристики: среднее время выполнения и количество посимвольных сравнений. Тестирование проводилось на трех наборах данных: \textit{случайных, обратных} и \textit{почти отсортированных} строках. Для каждого набора использовались подмассивы размеров от $100$ до $3000$ строк с шагом $100$.

\subsubsection{Время выполнения}

\begin{figure}[H]
    \centering
    \includegraphics[width=\textwidth]{/home/slava/comparison-of-sorting-algorithms/results/plots/standard_time_required_datasets.png}
    \caption{Время выполнения стандартных алгоритмов сортировки на трех наборах данных}
    \label{fig:standard_time_required_datasets}
\end{figure}

По времени выполнения \texttt{QUICKSORT} оказался быстрее \texttt{MERGESORT} на всех трех типах данных. Это связано с тем, что \texttt{QUICKSORT} сортирует массив преимущественно на месте, тогда как \texttt{MERGESORT} использует дополнительный буфер и выполняет больше операций перемещения элементов при слиянии.

На обратно отсортированных и почти отсортированных данных резкого ухудшения алгоритма \texttt{QUICKSORT} не наблюдается, поскольку в реализации опорный элемент выбирался из середины текущего фрагмента массива. Такой способ выбора опорного элемента снижает риск перехода к худшему случаю на упорядоченных данных.

\subsubsection{Количество посимвольных сравнений}

\begin{figure}[H]
    \centering
    \includegraphics[width=\textwidth]{/home/slava/comparison-of-sorting-algorithms/results/plots/standard_char_comparisons_required_datasets.png}
    \caption{Количество посимвольных сравнений стандартных алгоритмов сортировки на трех наборах данных}
    \label{fig:standard_char_required_datasets}
\end{figure}

По количеству посимвольных сравнений результат отличается от результата по времени: \texttt{QUICKSORT} выполняет заметно больше посимвольных сравнений, чем \texttt{MERGESORT}. Это объясняется тем, что в процессе разбиения массива элементы многократно сравниваются с опорными значениями. \texttt{MERGESORT} выполняет сравнения более регулярно при слиянии уже отсортированных частей, поэтому число посимвольных сравнений оказывается ниже.

Таким образом, меньшее число посимвольных сравнений не всегда означает меньшее время выполнения. Для \texttt{MERGESORT} существенную роль играют накладные расходы на дополнительную память и копирование строк при слиянии, поэтому по времени он уступает \texttt{QUICKSORT}, несмотря на меньшее количество посимвольных сравнений.

\subsubsection{Дополнительный набор строк с общим префиксом}

\begin{figure}[H]
    \centering
    \includegraphics[width=0.85\textwidth]{/home/slava/comparison-of-sorting-algorithms/results/plots/standard_char_comparisons_common_prefix.png}
    \caption{Количество посимвольных сравнений стандартных алгоритмов на строках с длинным общим префиксом}
    \label{fig:standard_char_common_prefix}
\end{figure}

Дополнительный набор строк с общим префиксом показывает ключевую особенность строковой сортировки. Оба стандартных алгоритма многократно заново сравнивают совпадающие начальные части строк. Поэтому при наличии длинного общего префикса количество посимвольных сравнений существенно возрастает.

Этот результат показывает, что для строк стандартная оценка сложности сортировки на основе сравнений недостаточна сама по себе. Необходимо учитывать не только количество сравнений элементов массива, но и стоимость каждого сравнения строк, зависящую от длины общего префикса.



\newpage
\section{Эмпирический анализ адаптированных алгоритмов}

Рассмотрены четыре специализированных алгоритма строковой сортировки: тернарный \texttt{STRING QUICKSORT}, алгоритм \texttt{LCP MERGESORT}, сортировка \texttt{MSD RADIX SORT} и ее версия с переключением на \texttt{STRING QUICKSORT} при размере фрагмента меньше мощности алфавита. Для сопоставимости также были включены стандартные \texttt{QUICKSORT} и \texttt{MERGESORT}. В стандартном \texttt{QUICKSORT} и тернарном \texttt{STRING QUICKSORT} использовался одинаковый способ выбора опорного элемента: элемент из середины текущего фрагмента массива.

Для \texttt{MSD RADIX SORT} значение счётчика посимвольных операций -- это не количество попарных сравнений строк, а количество обращений к символам при распределении строк по корзинам.

\subsection{Сравнение времени выполнения}

\begin{figure}[H]
    \centering
    \includegraphics[width=\textwidth]{/home/slava/comparison-of-sorting-algorithms/results/plots/all_algorithms_time_required_datasets.png}
    \caption{Время выполнения стандартных и адаптированных алгоритмов}
    \label{fig:all_time_required}
\end{figure}

По времени выполнения лучшим алгоритмом оказался тернарный \texttt{STRING QUICKSORT}. Он стабильно быстрее стандартного \texttt{QUICKSORT}, поскольку не выполняет полное лексикографическое сравнение строк при каждом разбиении, а работает с символом на текущей глубине.

Несмотря на уменьшение повторных сравнений префиксов, \texttt{LCP MERGESORT} имеет дополнительные накладные расходы на хранение и обновление LCP-информации при слиянии. Поэтому на строках без специально выраженных длинных общих префиксов эти расходы оказываются существенными.

\texttt{MSD RADIX SORT} тоже не всегда выигрывает по времени. Это происходит из-за высокой стоимости распределения по корзинам и рекурсивной обработки большого числа фрагментов. Вариант с переключением на \texttt{STRING QUICKSORT} работает быстрее чистого \texttt{MSD RADIX SORT}, так как уменьшает накладные расходы на малых подмассивах.

\subsection{Сравнение количества посимвольных сравнений}

\begin{figure}[H]
    \centering
    \includegraphics[width=\textwidth]{/home/slava/comparison-of-sorting-algorithms/results/plots/all_algorithms_char_comparisons_required_datasets.png}
    \caption{Количество посимвольных сравнений и операций}
    \label{fig:all_char_required}
\end{figure}

По количеству посимвольных сравнений адаптированные алгоритмы значительно превосходят стандартные. Особенно заметно отличие от стандартного \texttt{QUICKSORT}, который многократно сравнивает строки с опорными элементами и повторно просматривает одинаковые начальные части строк.

Минимальное число посимвольных операций показывают варианты \texttt{MSD RADIX SORT}. Это соответствует идее алгоритма: строки не сравниваются попарно, а распределяются по символам начиная со старшей позиции. \texttt{LCP MERGESORT} также уменьшает число повторных сравнений по сравнению со стандартным \texttt{MERGESORT}, так как использует уже известную длину общего префикса при слиянии.

\subsection{Дополнительный набор строк с общим префиксом}

\begin{figure}[H]
    \centering
    \includegraphics[width=0.9\textwidth]{/home/slava/comparison-of-sorting-algorithms/results/plots/all_algorithms_time_common_prefix.png}
    \caption{Время выполнения всех алгоритмов на строках с длинным общим префиксом}
    \label{fig:all_time_common_prefix}
\end{figure}

На строках с длинным общим префиксом преимущество тернарного \texttt{STRING QUICKSORT} становится наиболее заметным. Алгоритм эффективно обрабатывает группы строк с одинаковыми символами на текущей позиции и переходит к следующей позиции только внутри равной группы.

При этом \texttt{MSD RADIX SORT} и его вариант с переключением выполняются медленнее, несмотря на малое число посимвольных операций. Это показывает, что практическое время определяется не только числом обращений к символам, но и накладными расходами на создание и обработку корзин.

\begin{figure}[H]
    \centering
    \includegraphics[width=0.9\textwidth]{/home/slava/comparison-of-sorting-algorithms/results/plots/all_algorithms_char_comparisons_common_prefix.png}
    \caption{Количество посимвольных сравнений и операций на строках с длинным общим префиксом}
    \label{fig:all_char_common_prefix}
\end{figure}

На наборе с общим префиксом стандартные алгоритмы выполняют существенно больше посимвольных сравнений, так как при каждом сравнении заново просматривают совпадающую начальную часть строк. Адаптированные алгоритмы уменьшают этот эффект: \texttt{STRING QUICKSORT} группирует строки по текущему символу, \texttt{LCP MERGESORT} использует информацию о длине общего префикса, а \texttt{MSD RADIX SORT} распределяет строки по символам без попарного сравнения.

\subsection{Итоговое сравнение с базовым алгоритмом}

\begin{figure}[H]
    \centering
    \includegraphics[width=\textwidth]{/home/slava/comparison-of-sorting-algorithms/results/plots/speedup_time_vs_standard_quicksort_max_size.png}
    \caption{Ускорение алгоритмов относительно стандартного \texttt{QUICKSORT} при $n = 3000$}
    \label{fig:speedup_quicksort}
\end{figure}

Итоговое сравнение при $n = 3000$ показывает, что наибольшее ускорение относительно стандартного \texttt{QUICKSORT} даёт тернарный \texttt{STRING QUICKSORT}. На трех наборах данных он работает в несколько раз быстрее базового алгоритма, а на строках с общим префиксом его преимущество становится максимальным.

Вариант \texttt{MSD RADIX SORT} с переключением на \texttt{STRING QUICKSORT} быстрее базового \texttt{QUICKSORT} на трех наборах данных, но проигрывает на строках с общим префиксом из-за накладных расходов. Чистый \texttt{MSD RADIX SORT} показывает малое число посимвольных операций, но по времени уступает из-за высокой стоимости корзин. \texttt{LCP MERGESORT} уменьшает число сравнений, однако по времени не является лучшим из-за дополнительных операций с LCP-значениями.




\end{document}